\section{Related Work}
\label{sec:related}

\subsection{Tactile Representation Learning and Cross-Embodiment Retargeting}
Self-supervised representation learning has emerged as a promising paradigm for tactile sensing, adapting masked autoencoding, contrastive learning, and joint-embedding predictive architectures (JEPA) \citep{assran2023self,bardes2024revisiting} to visuo-tactile images and taxel arrays.
\citet{higuera2024sparsh} pre-train visuo-tactile encoders via MAE, DINO, and JEPA objectives on RGB elastomer images, while \citet{feng2025anytouch} learn static-dynamic representations across multiple visuo-tactile sensors using sensor-specific tokens.
However, these approaches predominantly operate in sensor-native image or taxel spaces, making cross-embodiment transfer between articulated dexterous hands and parallel grippers difficult.
On the geometric side, parametric human hand models such as MANO \citep{romero2017embodied} provide a canonical surface parameterization widely used for kinematic hand retargeting.
Unlike prior kinematic-only retargeting pipelines, our dual-route framework maps \emph{both} articulated hand kinematics with fingertip deformation arrays and static parallel-gripper tactile pads onto the same canonical MANO UV atlas under fixed, non-adaptive intensity scaling.

\subsection{Tactile and Embodied Manipulation Benchmarks}
Existing tactile evaluation suites largely focus on single-sensor classification, material recognition, or end-to-end policy success in simulation.
While informative for specific hardware, such benchmarks conflate encoder representation quality with downstream policy capacity, decoder architecture, or sensor-specific visual artifacts.
Table~\ref{tab:benchmark_comparison} contrasts \bench{} with representative tactile representation and manipulation benchmarks across embodiment coverage, canonicalization, validity masking, and evaluation protocol.
\bench{} is uniquely designed to isolate \emph{frozen} encoder latent quality across embodiments via unified small-decoder reconstruction (L0), decoupled proxy/physical/temporal probing (L1), and controlled cross-domain generalization (L2).

\begin{table*}[t]
\centering
\small
\setlength{\tabcolsep}{4.5pt}
\begin{tabular}{lccccccc}
\toprule
\textbf{Benchmark / Suite} & \textbf{Canonical} & \textbf{Dexterous +} & \textbf{Validity} & \textbf{Frozen} & \textbf{Unified} & \textbf{Temporal} & \textbf{Cross-Emb.} \\
 & \textbf{Surface} & \textbf{Gripper} & \textbf{Mask} & \textbf{Probes} & \textbf{Decoder} & \textbf{Dynamics} & \textbf{Transfer} \\
\midrule
Sparsh Suite \citep{higuera2024sparsh} & -- & Gripper only & -- & \checkmark & -- & Partial & -- \\
AnyTouch Eval \citep{feng2025anytouch} & -- & Gripper only & -- & \checkmark & -- & \checkmark & -- \\
Sensor-Native Recon. Baselines & -- & Single & -- & -- & -- & -- & -- \\
\midrule
\textbf{\bench{} (Ours)} & \textbf{MANO UV} & \textbf{\checkmark} & \textbf{\checkmark} & \textbf{\checkmark} & \textbf{\checkmark} & \textbf{\checkmark} & \textbf{\checkmark} \\
\bottomrule
\end{tabular}
\caption{Comparison of \bench{} with existing tactile representation evaluation protocols. \bench{} unifies articulated dexterous hands and parallel grippers onto a canonical MANO UV atlas with explicit observation validity masks and standardized frozen-encoder probing.}
\label{tab:benchmark_comparison}
\end{table*}
